\subsection{Nekov\'a\v r's third CM alternative and the cubic twists at thirteen}

Let $F$ be totally real and let $f$ be a parallel-weight-two CM
Hilbert eigenform.  Write $L^+$ and $L'$ for the real coefficient
field and its CM quadratic extension in Appendix B.6 of
\citet[Appendix~B.6]{Nekovar2012}.

\begin{lemma}[The split-coefficient condition is sufficient]
\label{lem:thirteen-split-coefficient}
Suppose that a rational prime $\ell$ splits completely in a
finite Galois extension containing $L'$.  Then the Galois-element
conditions (A1)--(A2) of that paper hold at every coefficient
prime above $\ell$.  In particular, after restriction to the
field trivializing the finite central character and the CM
quadratic extension, the two one-dimensional characters have
independent infinite images in the sense required by
Theorem B$'$(c2).
\end{lemma}
\begin{proof}
Proposition B.6.5(6), printed p.~664, applies because the real
coefficient completion is $\mathbb Q_\ell$ and its prime splits
in $L'/L^+$.  Proposition B.6.5(2) identifies the resulting
condition with the infinitude of $\psi_2(\ker\psi_1)$.
The proof on p.~665 is also explicit.  Dependence would force
powers of both characters to be powers of the cyclotomic
character.  Conjugacy exchanges them, so their exponents are
equal.  The cyclotomic exponent has odd weight; a potentially
crystalline character cannot then have the required half-integral
Hodge--Tate weight.  This proves the stated independence.
\end{proof}

This lemma removes character independence as an obstruction at
completely split coefficient primes.  It does not remove the
separate hypothesis $f\in S_2(\mathfrak n,1)$ in Theorems B and
B$'$ (printed p.~593).

For the exponent-thirteen factor, put
$L=\mathbb Q(\zeta_{13})$, $H=\langle3\rangle$, and
$K_0=L^H$.  Its maximal real subfield is
$F_0=\mathbb Q(\sqrt{13})$.
Let $B$ be the genus-two quotient used in the exponent-thirteen
calculation and let $\chi$ be the cubic Dirichlet character
with $\chi(2)=\zeta_3$.  The already established decomposition is
\[
 V_\ell P_{13}\simeq
 (V_\ell\operatorname{Jac}(B)\otimes\chi)
 \oplus
 (V_\ell\operatorname{Jac}(B)\otimes\chi^2).
\]
All finite coefficient extensions implicit in this display are
understood.  Over $F_0$, the real multiplication on
$\operatorname{Jac}(B)$ yields two-dimensional constituents
$W$ with determinant the cyclotomic character.  One may see this
directly from the polarization: the Rosati-fixed real
multiplication field acts self-adjointly, and the alternating
pairing on each constituent has cyclotomic similitude.
Consequently
\[
 \det(W\otimes\chi^j)=\chi_{\ell}\chi^{2j},
 \qquad j=1,2.
\]
The cubic character remains nontrivial on $G_{F_0}$, since its
real cubic field is disjoint from the quadratic field $F_0$.
Thus the corresponding Hilbert forms have nontrivial cubic
nebentypus.  Extending their coefficient fields, including
choosing completely split coefficient primes, does not change
this determinant character.

\begin{proposition}[Exact limit of this application]
\label{prop:thirteen-nebentypus-gate}
The split-coefficient lemma verifies the Galois-element
hypothesis for the exponent-thirteen cubic twists, but Theorems
B and B$'$ of Nekov\'a\v r (2012) do not apply directly to those
twists, because their nebentypus is nontrivial.
The elementary substitutions into Theorem A described below
do not repair this missing hypothesis.
\end{proposition}
\begin{proof}
The determinant calculation proves the first assertion.
Theorem A, printed p.~592, permits a nontrivial nebentypus
$\omega$, but requires an auxiliary CM quadratic extension
$K/F_0$ and a finite-order character $\theta$ with
\[
 \theta|_{\mathbb A_{F_0}^{\times}}\,\omega=1;
 \tag{*}
\]
its conclusion concerns the representation twisted by $\theta$
over $K$.  Taking $\theta=1$ leaves the desired representation
unchanged but violates (*) for the cubic twists.

More generally, start with the untwisted form $f_0$, replace it
by $f_0\otimes\chi^a$, and choose
$\theta=\chi^b\circ N_{K/F_0}$.  On the base ideles the norm is
squaring, so (*) is $\chi^{2(a+b)}=1$.  Since $\chi$ has order
three, this forces $a+b=0$ modulo three.  The resulting
representation over $K$ is untwisted, and hence is not either
of the two required cubic constituents.  Taking the auxiliary
CM field equal to the CM field of $f_0$ also violates the
explicit CM-field exclusion in Theorem A.

Finally, base change to the real cubic extension which kills
$\chi$ does not give a rank-zero untwisted problem: its cyclic
base-change $L$-function includes the untwisted
$\operatorname{Jac}(B)$ factor, and that variety already has
a non-torsion rational point.  In particular the entire
base-changed representation cannot have zero finite Selmer
group.  One would still have to isolate and control the
nontrivial cubic isotypic components; Theorem B applied to the
whole base change does not do this.
\end{proof}

This calculation concerns the precise scope of Nekov\'a\v r's
theorems, not a failure of the Tate--Shafarevich conclusion.
The following Euler-system argument treats the nontrivial-nebentypus
CM constituents directly and proves the needed vanishing for almost
all completely split rational primes.
